\documentclass[10pt]{article}
\usepackage[a4paper,margin=25mm]{geometry}
\usepackage{amsmath,amssymb,amsthm,mathrsfs,longtable,hyperref}
\emergencystretch=2em
\newtheorem{theorem}{Theorem}[section]
\newtheorem{lemma}[theorem]{Lemma}
\newtheorem{proposition}[theorem]{Proposition}
\newtheorem{corollary}[theorem]{Corollary}
\newcommand{\PSH}{\operatorname{PSH}}
\newcommand{\FS}{\operatorname{FS}}
\newcommand{\MA}{\operatorname{MA}}
\newcommand{\vol}{\operatorname{vol}}
\newcommand{\an}{\mathrm{an}}
\newcommand{\OO}{\mathcal O}
\newcommand{\E}{\mathcal E}
\newcommand{\dd}{dd^c}
\title{Raw model moments and discrepancy-controlled weighted limits\\
\large A partial research paper for the unchanged original P16.5}
\author{Research draft relative to the two explicit caller grants}
\date{6 October 2026}
\begin{document}
\maketitle
\begin{abstract}
We correct an energy identity used in the supplied nef-model bridge.
For the stated ordering of intersection factors its coefficients are
$n-j$, not $j+1$. An actual ample model on the blowup of
$\mathbf P^2\times\mathbf P^1$ at a point distinguishes the formulas.
The corrected identity still gives a dimension-only raw nef moment
estimate, including zero generic volumes. We then prove a compact
Hausdorff Radon-net squeezing theorem for arbitrary limiting support.
Its error is controlled by the two section-envelope discrepancies.
When those discrepancies vanish almost everywhere, it gives fixed
original continuous approximants with absolute $L^1$ error for bounded
relative tests. A separate lower-semicontinuous transfer shows precisely
how a reference discrepancy controls integrability of every relative
finite-energy test. The required discrepancy-nullness for general
big-class outputs is not proved. Nor are unbounded weighted
identification, arbitrary-target attainment, comparison or uniqueness.
The canonical result is therefore \texttt{candidate=false}.
\end{abstract}

\section{Retained domain, actual objects and the scope of the result}
The positive-volume domain sought throughout is every actual big rational
line bundle $L$ on every smooth integral projective variety $X$ of
dimension $n\geq1$ over every algebraically closed, trivially valued
characteristic-zero field $k$. Put $K=X^\an$, $\theta=c_1(L)$ and
$V=\vol_X(L)>0$. There is no countability, dimension, finite-generation,
toric, skeletal or target-support restriction on this sought domain.
A rational bundle is in $\operatorname{Pic}(X)\otimes\mathbf Q$.
Sections always use cleared tensor powers of the actual bundle.
Intersections may use numerical classes; section identities may not.

Only G0 and G1 in the supplied assignment are mathematical grants.
G0 supplies the complete stated ample theory on its retained
geometrically unibranch projective carriers and fields. We use it
only for actual ample bundles, or on a fixed finite model followed by
an ample perturbation limit. G1 supplies relative multiplier generation
and its one full generically finite FS net on the specified resolved
equivariant actual models. It is preserved and its role is audited below;
it does not grant any general nonample closure, energy or equation.
The supplied Complete151, Root194 and Root195 papers are unreviewed data.
We reconstruct the finite argument and the squeezing argument below.
The broader genuine foundations are explicit conditions of our
big-class applications, rather than silently accepted theorems.

The results proved here are: the corrected finite-model identities and
raw moment estimate, the elementary simultaneous full-system cofinality
argument relative to ordinary algebraic positive-product theory, the
general measure squeezing and integration-transfer lemmas, and the
numerically zero endpoint in Section~\ref{zero}. These do not constitute
the requested complete big-class equation theory.

\subsection{Original-carrier cone and the proposed intrinsic data}
A full section FS scalar on an actual $B$ is
\[
 \phi=m^{-1}\max_j(\log\|s_j\|_{\rm triv}+\lambda_j),
 \qquad 0\ne s_j\in H^0(X,mB).
\]
It may have common-zero poles. The numerical cone convention used in
the supplied programme consists of single all-point decreasing nets of
such scalars on actual $B_i$ with $c_1(B_i)\to\theta$, with a finite
generic limiting value. Denote it by $\PSH(\theta)$. This definition
does not automatically prove closure under a second decreasing limit.
All its elements are actual usc functions on $K$.

One useful operation does follow directly. If $\Omega-L$ is actual ample,
a tail of any defining net has $\Omega-B_i$ ample. Tensor with a
generating system of a cleared ample power, all with zero weights.
Its scalar maximum is exactly zero at every analytic point, because a
generator is a unit at the valuation-ring center. This retains the
incoming scalar in the actual $\Omega$ class. Appending generators at
weight $-N$ gives the scalar $\max(\phi_i,-N)$; the product index
$(i,N)$ decreases to the original endpoint and has no common zeros.
G0 therefore gives genuine ample membership on the same $K$.
It supplies divisorial determination and finite divisorial values for
these embedded functions, but supplies no equality of ample and big
energies.

Canonical section norms are at most one; a nonzero section has norm
one at the generic trivial valuation $v_0$. Consequently the supremum
of a defining stage, and of its endpoint, equals its generic value.
Integer scaling, constants and this common ample embedding follow
without a general nonample closure theorem.

The intended minimal reference is $V_\theta=P_\theta(0)$, normalized at
$v_0$. Its genuine membership is a foundation to prove. The intended
section-volume energy on a model obstacle $f_D$ is
\[
 \mathscr E(f_D)=\frac{\vol_Y(\pi_X^*L+D+cF)}{n+1}-cV
 \quad(c\text{ sufficiently large}),                         \tag{1}
\]
where $Y\to X\times\mathbf P^1$ is an actual smooth equivariant
modification, $F=\sum b_EE$ is the full central fiber, and
$f_D(v_E)=d_E/b_E$. Its extension is the genuine outer energy
\[
 E(u)=\inf_{\substack{f\in C(K)\\f\geq u}}\mathscr E(f),\qquad
 \E^1_\theta=\{u\in\PSH(\theta):E(u)>-\infty\}.                 \tag{2}
\]
In our big-class applications these formulas must have been proved
well defined, with $E(V_\theta)=0$, $E(u+c)=E(u)+cV$, and genuine
decreasing continuity. We do not assert those conclusions just by
writing (1)--(2).

For a positive Radon target $\nu$ of mass $V$ the proposed intrinsic
condition is
\[
 \nu\{V_\theta=-\infty\}=0,\quad
 D_\nu=\sup_{\substack{u\in\E^1_\theta\\u(v_0)=0}}
 \left\{\int(V_\theta-u)\,d\nu+E(u)\right\}<\infty.             \tag{3}
\]
The difference uses actual values off the reference poles.
It does not require absolute integrability of the reference.
Formula (3) is a specification of the equation programme; this paper
does not prove its sufficiency for an equation.

\section{A corrected raw finite-model identity}
Let $A$ be an actual nef rational bundle on a smooth projective $Z$,
$\dim Z=n\geq1$. A bounded relatively nef model scalar $\xi$
is represented by a vertical Cartier divisor on an actual equivariant
model. Write $T_\xi=(A,\xi)$ for its intersection slot, $T_0=(A,0)$,
$V_A=A^n$, and
\[
 e_A(\xi)=T_\xi^{n+1}/(n+1),\qquad
 \mu_A(\xi)=\MA_{A,\rm raw}(\xi),\qquad
 J_A(\xi)=\xi(v_0)V_A-e_A(\xi).
\]
Every model measure retains its central coefficient
$b_E(T_\xi|_E)^n$. Its mass is $V_A$.
These finite products and vertical Hodge inequalities are ordinary
model intersection calculus. Alternatively add a small actual ample
model to each slot, apply the finite ample calculus, and pass to the
finite polynomial limit. This operation does not extend an unbounded
nef equation.

For $n-1$ nef slots define
\[
 \|f\|_{T_1\cdots T_{n-1}}^2=-(0,f)^2T_1\cdots T_{n-1}\geq0.
                                                               \tag{4}
\]
Polarization gives Cauchy--Schwarz. All formulas are raw.

\begin{lemma}[Correct orientation of the tangent remainder]\label{identity}
Put $S=(0,\xi-\eta)$, so $T_\xi=T_\eta+S$. Then
\[
 T_\eta^{n+1}-T_\xi^{n+1}
       +(n+1)S T_\eta^n
 =-\sum_{j=0}^{n-1}(n-j)S^2T_\xi^jT_\eta^{n-1-j}.             \tag{5}
\]
In particular
\[
 J_A(\xi)=\frac1{n+1}\sum_{j=0}^{n-1}(n-j)
                  \|\xi\|_{T_\xi^jT_0^{n-1-j}}^2.             \tag{6}
\]
\end{lemma}
\begin{proof}
First telescope the difference of $(n+1)$st powers as
$S\sum_{k=0}^nT_\xi^kT_\eta^{n-k}$. Subtract
$(n+1)ST_\eta^n$. For each $k\geq1$ telescope
$T_\xi^kT_\eta^{n-k}-T_\eta^n$ again. The summand indexed by
$j$ occurs for $k=j+1,\ldots,n$, hence exactly $n-j$ times.
This proves (5). For $\eta=0$,
$(0,\xi)T_0^n=\xi(v_0)V_A$, since the trivial model measure is
$V_A\delta_{v_0}$. Divide by $n+1$ and use (4).
\end{proof}

With the factors in (5), coefficients $j+1$ are incorrect.
They occur in the supplied Root194 equation (6), and in the displayed
primary BJ v4 equation (3.22). Reading the PDF image, as well as its
rendered transcription, confirms that this is the printed orientation.
Reversing the index also requires reversing the two groups of factors.
We use (5), not the printed identity. The positive weights $1,\ldots,n$
are unchanged as a set, so the dimension-only comparability used in
the moment estimate survives.

\subsection{An actual ample model distinguishing the identities}
Take $Z=\mathbf P^2$, $A=\OO(1)$ and blow up
$(p,0)\in\mathbf P^2\times\mathbf P^1$. Let $E$ be the exceptional
prime, $H$ the pulled hyperplane and $0<a<1$ rational.
The actual model bundle $H-aE$ is relatively ample: on the strict
central component $\operatorname{Bl}_p\mathbf P^2$ its restriction is
the ample $H-aE_p$, on the exceptional component it is
$\OO_{\mathbf P^2}(a)$, and on other fibers it is $\OO(1)$.
All rational coefficients are cleared.
Since $E\simeq\mathbf P^2$, $E|_E=\OO(-1)$ and $H|_E$ is trivial,
\[
 H^3=H^2E=HE^2=0,\qquad E^3=1.
\]
For the scalar $\xi$ of $-aE$, its generic value is zero and
\[
 e_A(\xi)=-a^3/3,\quad J_A(\xi)=a^3/3,\quad
 \|\xi\|_{T_0}^2=0,\quad\|\xi\|_{T_\xi}^2=a^3.                \tag{7}
\]
Thus (6) gives $a^3/3$; the supplied $j+1$ formula gives $2a^3/3$.
The two central multiplicities are one. The raw model measure has
weights $1-a^2$ at $v_0$ and $a^2$ at the exceptional divisorial
valuation; its integral of $\xi$ is $-a^3$, as the derivative of
$-a^3/3$ requires. This is a mathematical finite-model
counterexample to that identity, not a counterexample to P16.5.

\section{The dimension-only nef moment estimate survives}
Set
\[
 d_A(\xi,\eta)=\max_{0\leq j<n}
       \|\xi-\eta\|_{T_\xi^jT_\eta^{n-1-j}}^2.
\]
The tangent remainder
$e_A(\eta)-e_A(\xi)+\int(\xi-\eta)\,d\mu_A(\eta)$ equals the
right side of (5) divided by $n+1$ with its signs reversed.
It is comparable to $d_A(\xi,\eta)$, with positive constants depending
only on $n$. It is convex in $\xi$ by concavity of the finite nef
energy. Also the expansion at the average gives
\[
 d_A(\xi,\eta)\asymp_n
 \|\xi-\eta\|_{T_{(\xi+\eta)/2}^{\,n-1}}^2.                   \tag{8}
\]

Here are the quantitative slot changes, to specify their scope.
Let $f=\xi-\xi'$, $\tau=(\xi+\xi')/2$,
$a=d_A(\xi,\xi')$, and
$b=\max(d_A(\xi,\eta),d_A(\xi',\eta))$.
Put $b_j=\|f\|_{T_\eta^jT_\tau^{n-1-j}}^2$.
Then $b_0\leq C_na$ by (8), and $b_{n-1}\leq4b$ by the
seminorm triangle inequality. In the difference $b_{j+1}-b_j$ replace
one $T_\tau$ by $T_\eta$ and split the second occurrence of $f$
into $T_\xi-T_{\xi'}$. Cauchy--Schwarz and
$T_\xi,T_{\xi'}\leq2T_\tau$ give
\[
 b_{j+1}-b_j\leq C_n\sqrt{b_j\,d_A(\tau,\eta)}
                    \leq C_n'\sqrt{b_jb}.                  \tag{9}
\]
The second inequality follows from convexity of the tangent remainder.
If $a\leq b$, induction yields
$b_j\leq C_na^{2^{-j}}b^{1-2^{-j}}$; if $a>b$, use
$b_{n-1}\leq4b$.
Applying this bound to two differences through a third potential gives,
with $d=d_A(\xi,\eta)$ and
$b=\max(d_A(\xi,\zeta),d_A(\eta,\zeta))$,
$d\leq C_n b^{2^{1-n}}\max(d,b)^{1-2^{1-n}}$.
If $d\geq b$, divide by $d^{1-2^{1-n}}$ and raise to the reciprocal
exponent; if $d<b$ the bound is immediate. This proves the
quasi-triangle bound for $d_A$.
For $n=1$ the empty product seminorm and its triangle inequality give
the same conclusion directly. Expanding the $(n-1)$st power of an
average of nef slots then gives
\[
 \|f\|_{T_{\eta_1}\cdots T_{\eta_{n-1}}}^2
 \leq C_n d_A(\xi,\xi')^{2^{1-n}}
 \max\{d_A(\xi,0),d_A(\xi',0),d_A(\eta_i,0)\}^{1-2^{1-n}}.
                                                               \tag{10}
\]
This is a finite algebraic argument on actual models. No unbounded
big-class regularization is used.

\begin{theorem}[Uniform raw finite moment]\label{moment}
For bounded relatively nef model scalars $h,\psi$ on the same actual
nef $A$, with $g_h=h(v_0)$, there is a dimension-only $C_n$ such that
\[
 \int(g_h-h)\,d\mu_A(\psi)
 \leq C_n J_A(h)^{\alpha}J_A(\psi)^{1/2}
         \max\{J_A(h),J_A(\psi)\}^{1/2-\alpha},
 \qquad\alpha=2^{-n}.                                     \tag{11}
\]
This includes $A^n=0$ and uses no volume denominator.
\end{theorem}
\begin{proof}
Shift both generic values to zero. The trivial model integral of $h$
is zero. The exact difference of the $n$th powers starts at $j=0$:
\[
 \int h\,d(\mu_A(\psi)-\mu_A(0))
 =\sum_{j=0}^{n-1}(0,h)(0,\psi)T_\psi^jT_0^{n-1-j}.
\]
Apply Cauchy--Schwarz to each term. The $\psi$ seminorm is at most
$d_A(\psi,0)^{1/2}$. The $h$ seminorm is bounded by the square root
of (10) with $\xi=h,\xi'=0$. Use (6) to replace $d_A$ by $J_A$
and sum the $n$ terms. Both normalized potentials are nonpositive,
so the desired left side has the correct sign.
Undo the constant shifts. Zero energies follow either directly from
the seminorm formulas or by finite polynomial continuity.
If $V_A=0$, positivity and the raw mass already force
$\mu_A(\psi)=0$; no probability normalization is possible.
\end{proof}

\section{Full moving systems: what is algebraic and what is conditional}
We retain ordinary characteristic-zero resolution, section-volume
asymptotics, and the algebraic positive-product theory of BFJ as
published background. These are not nonample analytic equation grants.
For actual big $N_1,\ldots,N_p$ on a smooth projective $d$-fold,
resolve all full degree-$m$ base ideals simultaneously and write
\[
 \rho_m^*N_i=C_{m,i}+G_{m,i}/m,\qquad
 mC_{m,i}\text{ generated},\quad G_{m,i}\geq0.
\]
Use factorial degrees clearing all denominators and nonempty systems.
Multiplication of full systems shows $C_{m\ell,i}\geq C_{m,i}$
on common dominations. These are actual divisor comparisons.

To check cofinality, fix finitely many nef competitors $\beta_i$
on one determination, with $N_i-\beta_i$ pseudoeffective.
Choose a small actual ample $Q_0$ with every $N_i-Q_0$ big.
For rational competitors and rational $\epsilon>0$ put
$Q_i=(1-\epsilon)\beta_i+\epsilon Q_0$; choose actual rational
representatives. Both $Q_i$ and $N_i-Q_i$ have the required
ample and big classes. In one sufficiently divisible degree $q$,
$qQ_i$ is generated and $q(N_i-Q_i)$ has an effective divisor $R_i$.
Multiplication by its section embeds the generated subseries in the
full original system, using proper birational normal section descent.
Thus $G_{q,i}\leq R_i$ and $C_{q,i}\geq Q_i$.
Nef-product monotonicity, followed by $\epsilon\downarrow0$ on the
fixed determination, proves cofinality in the defining least upper
bound. Real competitors follow by ample approximation.
BFJ's directed compactness then gives
\[
 C_{m,1}\cdots C_{m,p}\longrightarrow
                       \langle N_1\cdots N_p\rangle.         \tag{12}
\]
For a fixed signed Cartier test $H$ use a fixed decomposition into
two nef tests and subtract their limits. Do not pair an effective-order
comparison directly with a signed test.

For three model obstacles $0,D,H$ in a shared big $L$, choose one
sufficiently large time twist $cF$. After removing the same
horizontal full-$|mL|$ factor $T_m$, their complete moving bundles have
actual forms
\[
 C_{m,Z}=\rho_X^*A_m+Z+cF-G_{m,Z}^{\rm vert}/m,\quad
 C_{m,0}=\rho_X^*A_m+cF,\quad Z=0,D,H.                       \tag{13}
\]
The residual ideals are units off zero, including at infinity, because
the full generic moving system is free there. They define bounded
nef model scalars $\psi_{m,Z}$ on the same actual $A_m$.
The finite moment is
\[
 \int(-\psi_{m,H})\,d\mu_{A_m}(\psi_{m,D})
             =(C_{m,0}-C_{m,H})C_{m,D}^n.                   \tag{14}
\]
The factors $b_E$ in the measure and $b_E^{-1}$ in the scalar
cancel at each component, rather than being discarded.
By (12) the limit of (14) is the difference
\[
 \langle\alpha_0\alpha_D^n\rangle-
 \langle\alpha_H\alpha_D^n\rangle,\quad
 \alpha_Z=\pi_X^*L+Z+cF.                                   \tag{15}
\]
Positive products have not been made multilinear.

The further limits $A_m^n\to V$, $\psi_{m,Z}(v_0)\to0$, and
$e_{A_m}(\psi_{m,Z})\to E(P_\theta(f_Z))$ have different
logical origins. The first is algebraic volume approximation.
The other two use the genuine section-envelope and energy foundations.
When those foundations hold, (11) bounds (15) by
\[
 C_n B(u_H)^\alpha B(u_D)^{1/2}
          \max\{B(u_H),B(u_D)\}^{1/2-\alpha},\quad B(u)=-E(u).
                                                               \tag{16}
\]
At finite $m$ one must also retain
$-\psi_{m,H}(v_0)A_m^n$. The common generic class and the
generic error are essential.
Equating (15) to an original-carrier relative integral is still
a weighted-identification problem, addressed conditionally next.

\section{Squeezing at arbitrary Radon outputs}
The next arguments are measure theory on the original compact carrier.
They are independent of nonample equation foundations.

\begin{lemma}[Radon-net semicontinuity]\label{port}
If positive finite Radon measures $\nu_i\to\nu$ weakly on a compact
Hausdorff $K$, bounded lsc tests have the lower integral inequality
and bounded usc tests have the upper integral inequality, for nets.
\end{lemma}
\begin{proof}
A bounded lsc function is a directed supremum of continuous minorants.
Its Radon integral is their supremum: partition its range into finitely
many levels, choose compact subsets of the corresponding open
superlevels by inner regularity, and use continuous cutoffs on them.
Finite maxima combine the minorants. Weak convergence on each
minorant proves the lower inequality. Negation proves the upper one.
This uses no metrizability or subsequence extraction.
\end{proof}

\begin{theorem}[Quantitative arbitrary-support squeeze]\label{squeeze}
Let $\nu_i\to\nu$ be positive Radon measures with converging masses on
compact Hausdorff $K$. Let $|a_i|\leq C$ on their supports. Let
$U_q$ be increasing open sets with $\nu(G)=\nu(K)$ for
$G=\bigcup_qU_q$. On $U_q$ let $\ell_q,r_q\in[-C,C]$ be respectively
lsc and usc, with $\ell_q\leq a_i\leq r_q$ on the measure supports
there, eventually for each fixed $q$.
Suppose a Borel $a\in[-C,C]$ on $G$ satisfies
$\ell_q\leq a\leq r_q$, and the pointwise limits on $G$ are
$\ell_\infty,r_\infty$. Then
\[
 \int\ell_\infty\,d\nu
 \leq\liminf_i\int a_i\,d\nu_i
 \leq\limsup_i\int a_i\,d\nu_i
 \leq\int r_\infty\,d\nu.                                \tag{17}
\]
In particular, if $\ell_\infty=a=r_\infty$ almost everywhere, then
for every $\epsilon>0$ there is an actual $b\in C(K)$, $|b|\leq C$,
such that
\[
 \int|a-b|\,d\nu<\epsilon,\qquad
 \int|a_i-b|\,d\nu_i<\epsilon\quad\hbox{eventually}.          \tag{18}
\]
The degree $q$, compact collar and continuous $b$ are fixed before
the index $i$ tends to its limit.
\end{theorem}
\begin{proof}
For each $q$ choose compact $F_q\subset U_q$ with
$\nu(U_q\setminus F_q)<1/q$ and a continuous cutoff $0\leq\chi_q\leq1$
equal to one on $F_q$, supported in $U_q$. Then
$\int(1-\chi_q)d\nu\to0$.
The functions $\chi_q\ell_q,\chi_qr_q$, extended by zero, are
bounded lsc and usc on $K$. At the collar boundary their absolute
values are bounded by $C\chi_q$, which tends to zero.
Apply Lemma~\ref{port}, and bound the omitted part of each integral
by $C(\nu_i(K)-\int\chi_qd\nu_i)$. Send $i$ to its net limit
with $q$ fixed, then send $q$ to infinity.
Bounded convergence and the vanishing collar error prove (17).

For (18), a bounded Borel function has continuous $L^1(\nu)$
approximants on compact Hausdorff $K$: approximate by simple functions,
sandwich their Borel sets between compact and open sets by Radon
regularity, and apply continuous cutoffs. Clipping retains $|b|\leq C$.
Choose one such $b$ with arbitrarily small error. The nonnegative
function
\[
 R_{q,b}=\chi_q\max(r_q-b,b-\ell_q,0)
\]
extended by zero is bounded usc. On the supports in $U_q$ it
majorizes $\chi_q|a_i-b|$. Consequently
\[
 \limsup_i\int|a_i-b|d\nu_i
 \leq\int R_{q,b}d\nu+
              2C\left(\nu(K)-\int\chi_qd\nu\right).
\]
On $G$, the maximum is at most $|a-b|+(r_q-\ell_q)$.
Choose $q$ and its collar so both the width integral and collar
error are small, then fix $b$ with small limiting error.
This proves (18) with room in the strict inequalities.
If $C=0$, all statements are immediate.
\end{proof}

\section{The actual section quotient and its two defects}\label{defects}
Assume genuine model envelopes $u_H=P_\theta(f_H)$ and
$V_\theta=P_\theta(0)$, with full complete-system stages
$\phi_{q,H}\leq u_H$, $\phi_{q,0}\leq V_\theta$ increasing in
divisibility order. Their actual lattices imply
\[
 |\phi_{q,H}-\phi_{q,0}|\leq C
\]
where finite: shift each lattice by the two bounding time powers of
the bounded obstacle. Envelope monotonicity gives
$|u_H-V_\theta|\leq C$ where the reference is finite.
Take a cofinal factorial sequence of nonempty degrees.
Set $Q_H=\sup_q\phi_{q,H}$ and $Q_0=\sup_q\phi_{q,0}$.
We do \emph{not} assume their equality to the usc envelopes at
every original point.

Let $U_q$ be the open set where both degree-$q$ stages are finite.
On it define
\[
 \ell_q=\max(-C,\phi_{q,H}-V_\theta),\qquad
 r_q=\min(C,u_H-\phi_{q,0}).                                \tag{19}
\]
They are lsc and usc, respectively. For every sufficiently divisible
$m$, the changing quotient $a_m=\phi_{m,H}-\phi_{m,0}$
lies between them on $U_q$. On every pushed auxiliary divisorial atom
this quotient equals the actual $\psi_{m,H}$, by the common
horizontal-factor identity. No supported-point infinity is subtracted.

Suppose the original weak limiting output $\nu$ gives full mass to
$G=\bigcup_qU_q$. Put $a=u_H-V_\theta$ there and
\[
 \Delta_H=u_H-Q_H,\qquad\Delta_0=V_\theta-Q_0.
\]
These are nonnegative, possibly infinite, Borel discrepancies.
Theorem~\ref{squeeze} yields the explicit bound
\[
 \begin{split}
 \int a\,d\nu-\int\min(2C,\Delta_H)d\nu
 &\leq\liminf_m\int a_m\,d\nu_m,\\
 \limsup_m\int a_m\,d\nu_m
 &\leq\int a\,d\nu+\int\min(2C,\Delta_0)d\nu.
 \end{split}                                               \tag{20}
\]
Clipping in (19) only reduces these errors.
When both discrepancies vanish $\nu$-almost everywhere,
the general-support weighted identity and (18) follow.
These are absolute $L^1$ errors against actual continuous functions
on the original $K$ for the \emph{bounded relative quotient}.
They do not prove absolute integrability of $u_H$ or $V_\theta$,
or uniform tail control for an unbounded finite-energy test.

For finite divisorial outputs, the common actual ample embedding and
the ample separator imply no usc-supremum gap at each output atom.
For completeness, a separator $q_z\in\operatorname{PL}(K)$ with
$\sup(h-q_z)=(h-q_z)(z)$ for the common ample cone gives
$\sup_i h_i\leq\sup_i h_i(z)-q_z(z)+q_z$. Taking usc
regularization preserves this bound by continuity of $q_z$,
so evaluation at $z$ has no gap. This is the all-point application
of the primary finite-vertex argument, entirely in the ample
containing class. It identifies the conditional atomic moment (15)
with its original relative integral, recovering the finite-atom
mechanism in Root195. At general outputs, (20) displays precisely
the missing hypothesis instead of deleting it.

\subsection{A non-atomic bounded counterexample to an automatic extension}
On $K=[0,1]^2$ set $\phi_q(x,y)=\min(1,qx)$ and
let $\nu_m$ be Lebesgue mass one on the segment $x=m^{-2}$.
Then $\nu_m\to\nu$, Lebesgue mass one on $x=0$.
The increasing continuous stages have supremum
$Q(x,y)=1$ for $x>0$ and $Q(0,y)=0$; its usc regularization is
$u=1$ everywhere. With reference zero and $a_m=\phi_m$,
\[
 \int a_m\,d\nu_m=1/m\to0,\qquad\int u\,d\nu=1.             \tag{21}
\]
All tests are bounded by one; all first moments are uniformly bounded.
The defect $u-Q$ equals one on the entire limiting support.
This is a compact-space mathematical example, not a constructed
non-Archimedean output. It proves that weak convergence plus the
listed boundedness and stage monotonicity alone do not establish the
general weighted identification.

\section{How a reference defect controls every relative test}
\begin{theorem}[Lower-semicontinuous transfer]\label{integration}
Let $\nu_i\to\nu$ be positive Radon measures on compact Hausdorff $K$.
Let $\phi_q$ be increasing extended-continuous functions, finite on
increasing open $U_q$, with $\phi_q\leq V_*$ and $\sup_q\phi_q=Q_*$
on $G=\bigcup U_q$. Assume $\nu(G)=\nu(K)$ and $V_*$ is finite there.
Let an actual usc test $v\leq V_*$ satisfy
\[
 \int(V_*-v)d\nu_i\leq A(v)<\infty,
\]
with the relative differences defined on the supports.
Then
\[
 \int_G(Q_*-v)_+d\nu\leq A(v).                              \tag{22}
\]
If $R_\nu=\int_G(V_*-Q_*)d\nu<\infty$, then
\[
 \int(V_*-v)d\nu\leq A(v)+R_\nu.                            \tag{23}
\]
In particular discrepancy nullness gives the original moment bound.
\end{theorem}
\begin{proof}
On $U_q$, $(\phi_q-v)_+$ is nonnegative lsc, possibly infinite.
For a continuous $0\leq\chi\leq1$ supported in $U_q$,
extend $\chi\min(N,(\phi_q-v)_+)$ by zero. It is bounded lsc;
nonnegativity gives the lower inequality at a boundary point.
It is at most the relative gap on each measure support.
Lemma~\ref{port} therefore bounds its $\nu$ integral by $A(v)$.
First let $N\to\infty$, then exhaust $U_q$ in integral by compact
sets and cutoffs, to obtain
$\int_{U_q}(\phi_q-v)_+d\nu\leq A(v)$.
Extended by zero these nonnegative functions increase with $q$.
Monotone convergence proves (22).
Finite $R_\nu$ forces $Q_*$ finite almost everywhere.
There
$0\leq V_*-v\leq(Q_*-v)_++(V_*-Q_*)$,
including infinite test values. Integrate to get (23).
\end{proof}

In the big-class application, let $\nu_i$ be actual model-envelope
outputs of costs at most $B$, converging weakly to a target-independent
output $M(u)$. Conditional on the genuine foundations and the atomic
identification in Section~\ref{defects}, decreasing genuine
nonpositive model-majorant envelopes of a normalized
$v\in\E^1_\theta$ give the estimate
\[
 A(v)=C_n B(v)^\alpha B^{1/2}
                   \max\{B(v),B\}^{1/2-\alpha}.             \tag{24}
\]
Indeed apply (16) at a fixed model output first, then let the
nonpositive majorant envelopes decrease. Their generic values are
zero because they dominate normalized $v$ and lie below zero.
Their energies decrease to $E(v)$; monotone integration on a finite
output support gives (24). This avoids subtracting varying generic
constants from an arbitrary decreasing net.

If $M(u)$ gives full mass to $G$, (22) is thus available for every
actual test $v$. What is still needed to deduce the same estimate for
$V_\theta-v$ is reference discrepancy nullness, or at least
finite $R_{M(u)}$. Neither has been proved here. Algebraic pole-nullness
only controls the complement of $G$, not the discrepancy inside $G$.
Equations (22)--(24) expose that difference and do not assume it away.

\section{Audit of the supplied broader foundations}
The preceding conditional applications require the following exact
additional conclusions on the original carrier.
\begin{enumerate}
\item Genuine decreasing closure of the single-net cone, model envelopes,
and the identity
\[
 P_\theta(f)=(\sup\hbox{ full section minorants})^*.
\]
G1's one net alone is not this theorem. In the supplied relative-floor
argument the old pole comparison uses $h\leq\psi+C$ and an old floor
$F\geq(1-\delta)\psi-k$, giving
$h-F\leq\delta\psi+C+k$ near its poles. The assigned comparison value
there is $-\infty$, not a subtraction of infinities.
Finite compact constraints and the separate positive guard are needed
to flatten the incoming nets. No abstract twice-closure theorem is used
as a premise in this paper.
\item Actual normalized compactness, cofinal nonpositive model-majorant
envelopes, and the genuine energy in (1)--(2).
The original section lattice has weights $w_{mj}$ with
$h^0(Y,m(\pi_X^*L+D+cF))=\sum_j\max(w_{mj}+mc+1,0)$.
This accounts for the section-volume intercept. A bounded auxiliary
nef energy is not automatically the energy of the polar physical
section function $q_{T_m}+\psi_m$.
\item Every transfer from a moving subspace to the full original section
space retains the missing-rank term
$(\min f)(V-A_m^n)$, with $f$ a fixed original continuous majorant.
The section-degree limit for the fixed moving space precedes the
moving-degree limit. The generic-value error in (11) precedes the
genuine unbounded test limit. Neither error can be discarded at a
finite stage.
\item A target-independent outer-energy measure $M(u)$ of mass $V$,
finite-model compatibility, original continuous derivatives of both
signs, and decreasing weak continuity. The supplied time-tail method
uses deficient measures with total variation error at most $B/c$;
its order is fixed $c$, then the continuous variation parameter tending
to zero, then $c\to\infty$. Relative estimates cannot replace that
absolute continuous variation.
\item Nullness of algebraic section-zero sets for $M(u)$, when used to
put full mass on $G$. Even if the supplied truncated-section argument
is verified, it does not show either $\Delta_0=0$ or
$\Delta_H=0$ almost everywhere for these outputs.
\end{enumerate}
These are a precise audit of what the conditional bridge uses, not
additional grants or theorem certifications. We have not independently
completed every broader foundation from Complete151 in this paper.
The positive conclusions proved in Sections 2--7 do not rely on those
conclusions without the explicit conditional hypotheses.

\section{Why the equation still does not follow}
For the proposed target (3), the supplied contraction argument would give
coercivity and bounded-cost maximizing sequences. An actual compact
cluster is insufficient: the objective is $E(u)+I_\nu(u)$, where
$I_\nu(u)=\int(V_\theta-u)d\nu$. Weak upper semicontinuity of $E$
does not imply upper semicontinuity of this sum. The lower
semicontinuity bounds proved above have the other sign for attainment.

Needed are strong fixed-target integration on weak cost-bounded sets,
or a separately justified original-function strong-completeness
argument with no energy/integral jump. Only after attainment and
stationarity have been proved may both continuous derivative signs
identify $M(u)=\nu$. No output here is defined by solving an inverse
target problem.

There also remains the extension of all supporting planes to every
genuine finite-energy base, independence of simultaneous unbounded
weighted approximations, and strict separation. For comparison the
original Borel set $\{u<v\}$ need not be open; weak convergence does not
permit its indicator to pass through changing measures.
Even uniqueness of an eventual variational maximizer would not prove
uniqueness of every operator solution without the supporting-plane
and separation statements. Our new bounded squeeze does not prove
uniform integrable tails for unbounded tests. Even abstractly, bounded
first moments do not suffice: on $[0,1]$ put $w(x)=1/x$ for $x>0$,
$w(0)=+\infty$, and
$\nu_m=m^{-1}\delta_{1/m}+(1-m^{-1})\delta_1$.
Then $\nu_m\to\delta_1$, $\int w\,d\nu_m=2-1/m$, while the limiting
integral is one. For every fixed truncation level $N>1$, the tail
integral is one whenever $m>N$. This is a measure-theoretic
diagnostic, not a non-Archimedean finite-energy counterexample.

\section{The additional raw-zero range actually proved}\label{zero}
Here let the actual rational bundle $L$ be numerically trivial, on the
same smooth projective $X$ with $n\geq1$. This includes nontorsion
$\operatorname{Pic}^0$ bundles in the numerical varying-actual-bundle
cone convention. Reference is the canonical zero scalar.

\begin{proposition}
In that convention $\PSH(0)=\mathbf R$. Its raw zero-class operator
and raw energy are zero, and its only positive mass-zero target is
the zero measure. The equation is attained and unique modulo constants.
\end{proposition}
\begin{proof}
For every rational $\epsilon>0$, $L+\epsilon A$ is actual ample.
Its weight-zero generating system has identically zero scalar.
These constant scalars, with classes tending to zero, define zero
genuinely; shifts give every real constant.

Conversely normalize $u\in\PSH(0)$ at $v_0$ to zero. Integer
multiplication of its defining section systems gives $mu\in\PSH(0)$.
The common ample embedding from Section 1 puts every $mu$ in the
normalized genuine cone of one fixed actual ample bundle.
By G0 compactness, evaluation at a fixed divisorial $z$ has a finite
uniform lower bound $-C_z$ independent of $m$. Thus
$-C_z\leq mu(z)\leq0$ for all $m$, forcing $u(z)=0$.
G0 all-point divisorial determination on that ample cone gives $u=0$
on the entire original $K$.

The positive ample perturbations have raw mass
$\epsilon^n(A^n)\to0$. Each continuous integral is bounded by its
sup norm times this mass, so their unique limit is the zero
original-carrier Radon measure. On constants their raw energies are
$c\epsilon^n(A^n)\to0$. Hence the zero energy has zero constant law
and every constant is finite energy. A positive Radon measure of mass
zero is literally zero. The intrinsic relative dual energy is zero.
All solutions are the constants just classified, and therefore form
one class modulo constants. Comparison and the two zero derivative
signs are immediate. There is no division by zero.
\end{proof}

This is the whole numerically trivial range in the stated convention,
not the whole nonbig pseudoeffective cone. For example on
$\mathbf P^1\times\mathbf P^1$ the nonzero nef
$p_1^*\OO(1)$ has zero volume. It has zero and nonconstant bounded
pullback FS scalars. Adding $\epsilon\OO(1,1)$ gives positive raw
top products of mass $2\epsilon(1+\epsilon)\to0$ for both.
Thus their raw zero top products agree although their scalar difference
is not constant. This independently explains why the general
zero-volume uniqueness statement cannot simply retain the ample
normalization and cone.

\section{Precise canonical boundary and reproducibility}
The complete original P16.5 is \emph{not solved}. The retained
positive-volume target domain remains all actual big rational $L$ on
all smooth projective $X$ over all the stated fields.
The unresolved assertions are:
\begin{enumerate}
\item Independent completion/audit of the genuine closure, energy and
outer-output foundations specified in Section 8.
\item For every such output, nonalgebraic/negligible discrepancy-nullness,
in particular $\Delta_0=0$ almost everywhere and the corresponding
test discrepancies; reference integrability whenever an absolute
potential identification needs it.
\item All unbounded mixed moments, relative integrability of every
genuine $\E^1$ test without conditional reference defects, and
weighted identification independent of simultaneous approximations.
\item Actual original continuous $L^1$ approximants and uniform tails
for unbounded tests, with collars fixed before degree and measure limits.
\item Strong integration for every intrinsic target (3), actual
variational attainment without jumps, and the stationarity equation.
\item Entire supporting planes, original Borel strict comparison,
strict separation and uniqueness of every finite-energy operator
solution modulo constants.
\end{enumerate}
No other P16 family item is claimed solved. No nonzero nonbig
pseudoeffective equation classification is claimed.

The exact input bytes and decoded task/assignment are preserved in
\texttt{frozen/}; nine whole supplied units and inclusive assignment-LF
coordinates are in \texttt{supplied/index.json}. The local source
application ledger distinguishes grants, ordinary algebraic background,
new proofs, conditional inputs and diagnostic source readings.
Original renderer coordinates, PDF page labels and local decoded LF
coordinates are not conflated.
Run \texttt{python3 rerun.py}. It checks preservation, source-coordinate
coverage, result support order and finite exact polynomial identities,
then compiles this paper twice with the existing pdfLaTeX.
Commands, outputs and environment versions are in \texttt{logs/}.
Compilation and finite arithmetic do not certify mathematics.

Only this workspace, the user-supplied input and public primary sources
were used. No agents, other campaigns, credentials, account settings or
runtime changes were used. No installation or resource-heavy computation
was launched. No operator-supervised compute facility was exposed.

\appendix
\section{Primary source read scope}
BJ is Boucksom--Jonsson, \emph{Global pluripotential theory over a
trivially valued field}, arXiv:1801.08229v4, the PDF dated
24 March 2022. Pertinent rendered proof bodies retained are
Lemma 3.24, the finite concavity and slot-change arguments through
Corollary 3.34, and the ample divisorial separator and nonnegligibility
proofs. The equation (3.22) orientation is checked against PDF page
index 42, printed page 43, and corrected by our independent algebra.
The printed Theorem 3.31 right-hand side also repeats the first
potential; our quasi-triangle argument uses the two different
differences through the third potential. The diagnostic ample
Lemma 7.30 proof is retained; its displayed telescoping sum must
start at $j=0$. None of its unbounded ample regularization is
imported into a big class.

BFJ is Boucksom--Favre--Jonsson, \emph{Differentiability of volumes of
divisors and a problem of Teissier}, arXiv:math/0608260v2,
posted 4 September 2007. Relevant ordinary algebraic positive-product
and volume passages are preserved with exact renderer lines.
Its explicit characteristic-zero field scope is renderer line 147.
Its rendered HTML has an automatic date that is not treated as a new
source version. Hodge index, resolution, ample generation, section
asymptotics and Fujita remain stated ordinary background; no recursive
reading of their whole proof literature is claimed.

A targeted search also located Darvas--Xia--Zhang,
arXiv:2302.02541v3, and related transcendental/complex results.
Only the returned portions were read; no original-carrier arbitrary
Radon big-class NA equation is imported from a snippet or from those
sources. This is not an exhaustive bibliographical absence claim.
The ordinary direct PDF retrieval failed with an operation-not-permitted
error. Connector-rendered proof transcripts and their literal coordinates
are preserved; no downloaded original PDF hash is invented.

\begin{thebibliography}{9}
\bibitem{BJ} S. Boucksom and M. Jonsson,
\emph{Global pluripotential theory over a trivially valued field},
\href{https://arxiv.org/pdf/1801.08229v4}{arXiv:1801.08229v4}.
\bibitem{BFJ} S. Boucksom, C. Favre and M. Jonsson,
\emph{Differentiability of volumes of divisors and a problem of Teissier},
\href{https://arxiv.org/html/math/0608260v2}{arXiv:math/0608260v2}.
\bibitem{ADD} S. Boucksom and M. Jonsson,
\emph{Addendum to global pluripotential theory over a trivially valued field},
\href{https://arxiv.org/pdf/2206.07183v2}{arXiv:2206.07183v2}.
Only G1's exact supplied signature is a caller grant.
\bibitem{DXZ} T. Darvas, M. Xia and K. Zhang,
\emph{A transcendental approach to non-Archimedean metrics of
pseudoeffective classes},
\href{https://arxiv.org/pdf/2302.02541v3}{arXiv:2302.02541v3}.
Diagnostic literature reading only.
\end{thebibliography}
\end{document}
